%% ma: L10-B15-0096
%% lop: 10
%% chuong: 6
%% bai: 15
%% dang: 10.15.4
%% loai: TN4
%% muc: NB
%% dap_an: "C"
%% nguon: "Ảnh giáo viên gửi 25/11/2026 (lô 2): Bài tập trắc nghiệm – Toán 10 Bài 15 Hàm số (sự biến thiên, đồ thị), Phần I Câu 7"
%% kien_thuc: "Lớp 10 Bài 15 (sự đồng biến, nghịch biến của hàm số)"
%% trang_thai: chinh-thuc
\begin{ex}
Trong các hàm số sau, hàm số nào giảm trên khoảng $(0;1)$?
\choice{$y=x^2$}{$y=x^3$}{\True $y=\dfrac1x$}{$y=\sqrt x$}
\loigiai{Với $0<x_1<x_2<1$ thì $\dfrac1{x_1}>\dfrac1{x_2}$ nên $y=\dfrac1x$ giảm; ba hàm còn lại tăng trên $(0;1)$.}
\end{ex}
